\chapter{Conclusion and Future Work}
\label{Conclusion_and_Future_Work}

\section{Summary of the Work}

This dissertation presented a digital watermarking framework that works across four carrier types: $256 \times 256$ BMP images, WAV audio, plain-text files, and video containers. The main contribution is a row-based least-significant-bit method for images. In this method, the MD5 hash of the user password is used to select the image row and the starting column where the watermark will be stored. This allows multiple independent watermarks to exist in the same image, each linked to a separate password. A sentinel character at column~0 of an occupied row prevents accidental overwrite, a 3--3--2 RGB split keeps the visual change small, and a self-delimited payload allows extraction without a separate length field.

The audio, text and video modules reuse the same payload format, the same password input and the same in-window log under a single Swing GUI. Audio uses single-bit LSB substitution on raw PCM samples preceded by an XOR encryption; text appends an XOR-encrypted block inside an HTML-style comment; video extracts the first frame, applies the image algorithm, and re-encodes losslessly with FFV1 so that the LSBs survive.

Experimental evaluation on standard test images, WAV clips, text files and short videos confirms five principal conclusions:

\begin{itemize}
    \item \textbf{Imperceptibility}: PSNR $\ge 61$~dB on $256 \times 256$ images, no audible difference on PCM audio, and identical rendering of stego-text in plain-text, Markdown and HTML viewers.
    \item \textbf{Multi-watermark capacity}: ten independent watermarks were successfully embedded into and recovered from the same \texttt{lena.bmp} cover under ten different passwords, validating hypothesis~\textbf{H1}.
    \item \textbf{Collision detection}: the row-occupation marker reliably refuses an eleventh embedding under a colliding password, validating hypothesis~\textbf{H2}.
    \item \textbf{Run-time}: the image, audio and text modules complete an embedding or extraction in less than $130$~ms on commodity hardware; the video module runs in a few seconds and is dominated by FFmpeg I/O.
    \item \textbf{Multi-modal parser reuse}: the same self-delimited timestamped payload is recovered by a single parser regardless of carrier, validating hypothesis~\textbf{H3}.
\end{itemize}

The qualitative comparison of Chapter~\ref{Experimental_results} positions the system as a multi-watermark, multi-modal alternative to single-modality LSB tools rather than a competitor to robust transform-domain watermarking. The system has been delivered as a single self-contained Swing application of approximately $1{,}500$ lines of Java code, accompanied by a comprehensive in-window algorithmic log that prints every internal step. This makes the project useful both as a working ownership-proof tool for lossless distribution scenarios and as a teaching reference for undergraduate courses on multimedia security.

One clear outcome of this work is that integration itself can be an important contribution. The project does not depend on one highly complex algorithm. Instead, it shows that a careful combination of simple ideas can still produce a useful and testable framework. That is an important lesson in software-oriented research, where system design often matters as much as the individual techniques inside it.

\section{Limitations}

In the spirit of a frank technical evaluation, the current prototype has the following limitations.

\begin{enumerate}[label=\textbf{L\arabic*.}]
    \item \textbf{Fragility under practical media transformations.} The current design is based mainly on spatial-domain LSB substitution. As a result, the embedded information is vulnerable to JPEG compression, audio transcoding, video recompression, resizing, cropping, and format conversion. This limits the framework to controlled or lossless distribution settings rather than open, real-world content pipelines.
    \item \textbf{Restricted generality of the image pipeline.} The image module is tightly coupled to a $256 \times 256$ BMP cover and to a row index in $[0,255]$. This makes the present implementation easy to explain, but it also limits scalability to arbitrary resolutions, modern image formats, and heterogeneous multimedia datasets.
    \item \textbf{Limited cryptographic assurance.} The current prototype uses MD5 for deterministic location derivation and XOR-based obfuscation in the audio and text modules. These choices are acceptable for a classroom prototype, but they do not provide modern guarantees of collision resistance, confidentiality, integrity, or message authentication in an adversarial setting.
    \item \textbf{Modest empirical scale.} The experimental corpus is small and covers only a narrow range of images, short WAV clips, simple text files, and short videos. The present evaluation therefore demonstrates feasibility, but it is not yet broad enough to claim strong generalization across diverse content types, resolutions, languages, or production-grade media.
    \item \textbf{Limited adversarial evaluation.} The system has not yet been tested against standard steganalysis methods, forensic re-processing, deliberate tampering, or hostile attempts to estimate or remove the embedded payload. This means the framework's security properties are presently measured more by successful recovery in benign settings than by resilience against an active adversary.
    \item \textbf{Video pipeline maturity.} The video module currently embeds only in the first frame and depends on external FFmpeg-based extraction and lossless FFV1 re-encoding. While this is functionally correct, it does not yet provide temporal redundancy, robustness across frame loss, or an efficient workflow for large-scale video deployment.
    \item \textbf{Prototype-level workflow support.} The system is presently optimized for demonstration through a Swing GUI. It does not yet offer a command-line interface, batch processing, automated regression tests, or reproducible benchmark scripts at the level expected from a mature research tool or deployable software package.
\end{enumerate}

\section{Future Work}

The following extensions would address the limitations listed above and would significantly broaden the practical applicability of the system.

\begin{enumerate}[label=\textbf{F\arabic*.}]
    \item \textbf{Robust watermarking under compression and geometric distortion.} A major next step is to design a hybrid embedding mode that combines the present spatial LSB method with transform-domain watermarking in DCT, DWT, or SVD coefficients. This would allow the system to retain at least partial recoverability after compression, scaling, cropping, and moderate signal processing.
    \item \textbf{Cryptographic hardening and authenticated payload protection.} The present MD5/XOR pipeline should be replaced with a modern key-derivation and encryption design such as HMAC-SHA256 or HKDF for deterministic indexing and AES-GCM for confidentiality and integrity~\cite{daemen2002design}. This would move the framework from a didactic prototype toward a security-aware system.
    \item \textbf{Generalization to arbitrary cover sizes and formats.} The row-based image method should be extended to general $W \times H$ images and modern lossless formats such as PNG and TIFF. More broadly, the framework should support modality-aware capacity estimation so that payload allocation is adapted to the actual size and structure of the input cover.
    \item \textbf{Steganalysis-aware embedding.} Future work should evaluate the framework against classical and modern detectors, including RS analysis, chi-square tests, and machine-learning-based steganalysis. Based on those results, the embedding rule can be upgraded toward lower-detectability alternatives such as LSB matching, matrix encoding, or adaptive region selection.
    \item \textbf{Temporal extension of the video module.} Instead of embedding only in frame~0, future versions should distribute the watermark across multiple frames or key frames. A temporally distributed design would reduce per-frame distortion, increase resilience to frame loss or editing, and better align the video module with practical multimedia workflows.
    \item \textbf{Larger-scale benchmarking and statistical validation.} A more mature evaluation should use substantially larger corpora across image, audio, text, and video modalities, together with repeated trials, comparative baselines, and statistical reporting. This would allow stronger claims regarding capacity, invisibility, recovery accuracy, and computational cost.
    \item \textbf{Tooling for deployment and reproducible experimentation.} The framework would benefit from a command-line interface, batch-processing mode, automated test suite, and benchmark harness. These additions would make the system easier to use in scripted workflows and would improve reproducibility for future academic work.
    \item \textbf{Ownership verification and forensic workflows.} An important longer-term extension is to connect the embedding system with practical verification workflows, including recipient-specific watermark assignment, audit logging, watermark provenance tracking, and post-distribution forensic recovery. This would strengthen the system's relevance beyond proof-of-concept laboratory settings.
\end{enumerate}

\section{Closing Remarks}

This project shows that a simple and understandable design can still solve a useful problem. The row-based image method, the overwrite check, and the shared payload format together make the framework practical for study and future extension. The main lesson from this work is that the way the components are combined matters a lot. The same architecture can later be improved with stronger cryptography, support for more file sizes, and better robustness against lossy compression.

At the same time, the project also makes its own limits clear. It works best in controlled conditions, and it should not be described as a complete answer to real-world copyright protection. Still, that does not reduce its value. As an academic project, it demonstrates clear design thinking, a working implementation, cross-media consistency, and a realistic path for future improvement.

For that reason, the work can be judged successful on its intended scale. It delivers a usable prototype, explains its methods in a transparent way, and leaves behind a structure that later work can strengthen rather than replace.
